\section{The Integral}

\subsection{Antiderivatives}
\begin{definition}
$F$ is an antiderivative of $f$ on $I$ if $F'(x) = f(x)$ for every $x \in I$.
Two antiderivatives of the same function differ by a constant, so
\[ \int f(x) \dx = F(x) + C. \]
\end{definition}

\subsection{Riemann sums and the definite integral}
Partition $[a,b]$ into $n$ subintervals of width $\Delta x_i$ and pick a sample point
$x_i^*$ in each. The Riemann sum is $\sum_{i=1}^n f(x_i^*) \Delta x_i$.

\begin{definition}[Definite integral]
If the Riemann sums converge to a single value as the maximum width tends to $0$
(independent of the partition and the sample points), $f$ is \emph{Riemann integrable} on
$[a,b]$ and we write
\[ \int_a^b f(x) \dx = \lim_{\max \Delta x_i \to 0} \sum_{i=1}^n f(x_i^*) \Delta x_i. \]
\end{definition}

Continuous functions on $[a,b]$ are integrable; so are bounded functions with at most
finitely many discontinuities.

\subsection{Properties}
\[
\begin{aligned}
&\int_a^a f \dx = 0, \qquad
\int_a^b f \dx = -\int_b^a f \dx, \\
&\int_a^b (\alpha f + \beta g) \dx = \alpha \int_a^b f \dx + \beta \int_a^b g \dx, \\
&\int_a^c f \dx + \int_c^b f \dx = \int_a^b f \dx, \\
&f \le g \text{ on } [a,b] \;\Longrightarrow\; \int_a^b f \dx \le \int_a^b g \dx.
\end{aligned}
\]

\subsection{Fundamental Theorem of Calculus}
\begin{theorem}[FTC, Part I]
If $f$ is continuous on $[a,b]$ and $F(x) = \int_a^x f(t) \dt$, then $F$ is differentiable on
$(a,b)$ and $F'(x) = f(x)$.
\end{theorem}

\begin{theorem}[FTC, Part II / Evaluation theorem]
If $f$ is continuous on $[a,b]$ and $F$ is any antiderivative of $f$ on $[a,b]$, then
\[ \int_a^b f(x) \dx = F(b) - F(a). \]
\end{theorem}

\begin{intuition}
Differentiation and integration are inverse operations: differentiate the integral and you get
$f$ back; integrate the derivative and you recover $F$ up to its endpoint values.
\end{intuition}

\subsection{Leibniz rule (variable limits)}
\[ \frac{d}{dx}\!\int_{a(x)}^{b(x)} f(t) \dt = f(b(x)) b'(x) - f(a(x)) a'(x). \]

\subsection{Mean Value Theorem for Integrals}
\begin{theorem}
If $f$ is continuous on $[a,b]$, then there exists $c \in [a,b]$ with
\[ f(c) = \frac{1}{b-a} \int_a^b f(x) \dx. \]
\end{theorem}
The right-hand side is the \emph{average value} of $f$ on $[a,b]$.

\subsection{Standard antiderivatives}
\begin{keybox}[Table to memorize]
\[
\renewcommand{\arraystretch}{1.3}
\begin{array}{l|l}
f(x) & \int f(x) \dx \\\hline
x^n\ (n \neq -1) & \dfrac{x^{n+1}}{n+1} + C \\[4pt]
\dfrac{1}{x} & \ln|x| + C \\[4pt]
e^x & e^x + C \\
a^x & \dfrac{a^x}{\ln a} + C \\[4pt]
\sin x & -\cos x + C \\
\cos x & \sin x + C \\
\sec^2 x & \tan x + C \\
\sec x \tan x & \sec x + C \\[2pt]
\dfrac{1}{1+x^2} & \arctan x + C \\[4pt]
\dfrac{1}{\sqrt{1-x^2}} & \arcsin x + C \\[4pt]
\dfrac{1}{x\sqrt{x^2-1}} & \arcsec|x| + C
\end{array}
\]
\end{keybox}
